% Vietnamese translation for OI Public Library
% Source-File: IOI中国国家候选队论文/2003/陆可昱/陆可昱.ppt
\documentclass[11pt]{article}
\usepackage{oipl-vn}

\title{Bản trình chiếu: Hợp thể tích các hình hộp chữ nhật}
\author{Lu Keyu}
\date{2003}

\begin{document}
\maketitle

\origtitle{Union volume of cuboids}
\sourcepath{IOI China candidate team papers / 2003 / Lu Keyu / Lu Keyu.ppt}

\section{Từ hợp diện tích hình chữ nhật}

Bài trình bày bắt đầu bằng bài quen thuộc: tính diện tích hợp của \(n\) hình
chữ nhật trên mặt phẳng. Cách làm gồm ba bước:
\begin{itemize}
  \item rời rạc hóa tọa độ;
  \item quét mặt phẳng theo từng dải;
  \item dùng cây đoạn để lưu trạng thái phủ trên tiết diện một chiều.
\end{itemize}

Điểm rời rạc là giao điểm của các cạnh hình chữ nhật, hoặc đường kéo dài của
chúng, với trục tọa độ. Đơn vị rời rạc là khoảng cách giữa hai điểm rời rạc
liền kề sau khi sắp thứ tự. Trong ví dụ ở slide, các điểm \(A,B,C,D\) trên
một trục có thể được đánh số \(1,2,3,4\).

Khi mặt phẳng được chia thành các dải bởi các đường quét, mỗi dải chỉ còn
một bài toán một chiều. Tiết diện của một dải có thể thu được từ tiết diện
của dải kề bằng một sửa đổi nhỏ, ví dụ thêm hoặc bớt một đoạn.

\section{Cây đoạn}

Cây đoạn trong slide là một cây nhị phân có gốc. Mỗi nút biểu diễn một
khoảng \([a,b]\). Nếu \(b-a>1\), đặt
\[
  c=\left\lfloor\frac{a+b}{2}\right\rfloor,
\]
rồi tách thành hai nút con \([a,c]\) và \([c,b]\). Với cách đánh số tĩnh,
gốc có nhãn \(1\); nút không lá \(i\) có con trái \(2i\) và con phải
\(2i+1\).

Trong bài diện tích hợp hai chiều, cây đoạn lưu độ dài hợp của các đoạn đang
được phủ trên tiết diện. Khi đường quét đi qua biên trái hoặc biên phải của
một hình chữ nhật, ta thêm hoặc xóa một đoạn và cập nhật lại độ dài phủ.

\section{Chuyển sang ba chiều}

Bài toán chính là tính thể tích hợp của \(n\) hình hộp chữ nhật trong không
gian ba chiều. Tư tưởng vẫn là rời rạc hóa và quét, nhưng lần này khi quét
theo một trục, tiết diện cần lưu không còn là đoạn một chiều mà là một mặt
phẳng gồm các hình chữ nhật.

Vì vậy slide thay cây đoạn một chiều bằng cấu trúc ``cây nhị phân hai tầng''.
Ta xây một cây đoạn trên trục \(x\) và một cây đoạn trên trục \(y\). Một cặp
nút \((x_1,y_1)\) biểu diễn một hình chữ nhật rời rạc trong mặt phẳng. Mảng
\[
  T[x_1][y_1]
\]
lưu thông tin của vùng có thành phần ngang là khoảng của nút \(x_1\) trên
cây \(x\), và thành phần dọc là khoảng của nút \(y_1\) trên cây \(y\).

\section{Thêm và xóa hình chữ nhật trên tiết diện}

Mỗi phần tử \(T[x_1][y_1]\) có trường \(C\), ghi số lần vùng tương ứng được
phủ hoàn toàn. Khi thêm hoặc xóa một hình chữ nhật trên tiết diện, ta chèn
khoảng ngang của nó vào cây \(x\) và thu được tập các nút cuối \(A_x\); đồng
thời chèn khoảng dọc vào cây \(y\) và thu được tập \(A_y\). Với mọi
\[
  p\in A_x,\qquad q\in A_y,
\]
ta sửa \(T[p][q].C\).

Slide có ví dụ chèn hình chữ nhật
\[
  [1,5]\times[3,7].
\]
Khoảng \([1,5]\) trên cây \(x\) đi đến các nút \(5,9,12\), còn khoảng
\([3,7]\) trên cây \(y\) đi đến các nút \(6,11,14\). Vì thế các nút có
trường \(C\) cần sửa là mọi cặp trong tích Descartes của hai tập này, chẳng
hạn \(T[9,14]\), \(T[9,6]\), \(T[5,11]\), \(T[12,6]\), v.v.

\section{Tính diện tích của mặt phẳng đang phủ}

Mỗi \(T[x_1][y_1]\) còn có trường \(M\), lưu diện tích hợp của các hình chữ
nhật trong vùng đó. Nếu
\[
  T[x_1][y_1].C>0,
\]
toàn bộ vùng được phủ và \(M\) bằng diện tích hình chữ nhật của nút. Nếu
\[
  T[x_1][y_1].C=0,
\]
thì \(M\) được cộng từ các vùng con.

Slide minh họa trường hợp vùng \(ABCD\) chưa phải lá trên cả hai cây. Gọi
\(E,F,G,H\) là trung điểm của bốn cạnh và \(I\) là giao điểm của hai đường
trung bình. Diện tích của \(ABCD\) được tách thành bốn phần
\[
  AEIH,\quad EBFI,\quad IFCG,\quad HIGD.
\]
Khi cộng một phần nhỏ như \(AEIH\), cần chú ý các vùng lớn hơn như \(ABFH\)
hoặc \(AEGD\) có thể đã được phủ. Nếu một trong hai trục đã ở nút lá, phép
cộng đơn giản hơn; slide chỉ nêu ý tưởng và không triển khai chi tiết từng
trường hợp biên.

\section{Cập nhật trường diện tích}

Khi thêm một hình chữ nhật, ngoài các nút cuối \(A_x,A_y\) dùng để sửa
\(C\), ta còn ghi lại mọi nút đã đi qua trên hai cây:
\[
  S_x,\qquad S_y.
\]
Với ví dụ \([1,5]\times[3,7]\), slide ghi
\[
  S_x=\{1,2,3,4,5,6,9,12\},
\]
\[
  S_y=\{1,2,3,5,6,7,11,14\}.
\]
Các giá trị \(M\) cần cập nhật nằm trong các cặp nút của hai tập này. Việc
cập nhật được thực hiện từ nút sâu hơn lên trên. Nhờ cách đánh số tĩnh của
cây nhị phân, nút sâu thường có nhãn lớn hơn, nên thứ tự cập nhật có thể
được xử lý thuận tiện bằng nhãn.

\section{Độ phức tạp và mở rộng}

Trong mỗi lần thêm hoặc xóa một hình chữ nhật trên mặt quét, số nút cuối trên
mỗi cây là cấp \(\log n\). Do đó sửa \(C\) tốn
\[
  O(\log^2 n).
\]
Số nút đi qua khi sửa \(M\) cũng là cấp \(\log n\) trên mỗi trục, nên sửa
\(M\) cũng tốn
\[
  O(\log^2 n).
\]

Khi tính hợp thể tích của \(n\) hình hộp, đường quét gặp \(2n\) mặt vào/ra,
nên tổng độ phức tạp được slide ghi là
\[
  O(n\log^2 n).
\]

Phần mở rộng của slide nêu rằng cùng tư tưởng có thể áp dụng cho không gian
\(d\) chiều: rời rạc hóa, quét, lưu khối tiết diện, và thay cây nhị phân hai
tầng bằng cây nhị phân \(d\) tầng. Tuy nhiên, số chiều càng lớn thì độ phức
tạp cài đặt càng tăng nhanh.

\end{document}
